\chapter{模型评价、选择与泛化}
\label{chap:evaluation}
\chaptergoal{能区分学习参数、选择方案和最终检查的用途；能从具体预测记录计算分类与回归指标；能解释基线、欠拟合、过拟合和交叉验证；能把数据处理与模型放在同一验证流程中。}

第3章说明参数怎样改变、训练损失怎样下降。现在需要回答另一个问题：学到的规则遇到新记录时是否仍然有用？把训练记录全部记住并不等于能处理新的对象。本章先讨论怎样留出用于检查的记录，再讨论怎样把预测好坏写成数值，最后将评价用于模型选择。

本章继续使用第2章介绍的葡萄酒数据：每条记录是一份酒样的13项分析指标，目标为三类来源。选择熟悉的数据，是为了把注意力放在评价过程，而不是同时适应新的研究对象。住房回归只作为已介绍的连续目标任务进行对照。

\section{训练、验证与独立测试}
\subsection{三种用途，不是三种文件格式}
\textbf{训练集（training set）}用于确定模型参数或其他预测规则。\textbf{验证集（validation set）}用于比较候选算法、预处理方式与参数配置。\textbf{测试集（test set）}是在方案确定后，用来检查新样本表现的数据。前两者共同组成开发过程；“独立测试”强调测试信息没有反过来参与方案选择。

可以把开发过程想成先做练习，再用一套模拟题选择学习策略，最后参加不用于改策略的考试。训练集对应用于形成规则的练习，验证集帮助选择规则，测试集用于最终检查。关键是信息怎样流动，而不是变量名怎样命名：反复根据某个所谓测试集改模型，它实际上已经承担验证用途。

\subsection{一个有记录数的划分例子}
假设有100条适合按独立样本随机划分的记录，可以先留20条作最终测试，余下80条用于开发。再把开发数据分成60条训练、20条验证。训练三个候选方案时，每个方案都只用相同的60条记录建立规则，再在相同20条验证记录上比较。选定方案后，可以按预先约定在全部80条开发记录上重新拟合，最后在最初保留的20条上评价。

60、20、20只是解释角色的人工例子，不是必须采用的比例。实际葡萄酒数据有178条，以下代码先留约20\%，得到142条开发记录和36条最终保留记录；开发部分内部再用交叉验证，详见4.5节。
\begin{lstlisting}[caption={建立后续章节共用的开发与保留划分}]
from sklearn.datasets import load_wine
from sklearn.model_selection import train_test_split
wine = load_wine(as_frame=True)
X_dev, X_test, y_dev, y_test = train_test_split(
    wine.data, wine.target, test_size=0.20,
    random_state=42, stratify=wine.target
)
print(len(X_dev), len(X_test))
\end{lstlisting}

\textbf{分层划分}尽量保持类别比例，避免偶然把某个少数类几乎全部分到一边。\textbf{分组划分}则按对象或实验组分配记录，例如同一天体的重复观测进入同一部分。\textbf{时间划分}按先后顺序检查对未来的预测。三者解决的不是同一个问题，选择应对应预期应用，而非机械地一律随机拆行。
\fig{ch04_split.pdf}{开发数据、内部验证与最终测试的角色。最终测试不是另一次用来挑选最好参数的验证。}

\subsection{一套固定口径贯穿后续实验}
本书新增的标准实验统一先按种子42保留20\%作为最终测试；剩余数据称为开发集。算法比较主要采用开发集上的五折分层交叉验证，折划分种子为7。为绘制训练曲线或追踪一个样本，某些实验再把开发集按种子7划分为75\%训练和25\%验证。两种实验分别标为“五折比较”和“固定验证示范”，不能把它们的分数当作同一条件的重复测量。

对鸢尾花，固定示范有90条训练、30条验证，另30条最终保留；对葡萄酒，有106条训练、36条验证，另36条最终保留。小数比例经程序取整，所以记录数比单纯报告百分比更准确。拆分索引保存在随书CSV中，之后各章无需重新随机寻找一个更好看的划分。

在初学阶段，保留测试数据的价值不只是“获得一个最终分数”，还在于让模型选择和评价的职责分开。若课程已经围绕某个标准数据反复教学多年，它不再是研究意义上完全陌生的新数据，但仍可用来练习正确流程。真实科学报告需要另外考虑数据选择历史，不能把课堂保留方式当作自动消除所有研究偏差的办法。

\subsection{按对象划分与按行划分回答不同问题}
随书构造了一个可以精确检查的反例。90个模拟对象各有5次测量，每个对象具有一个随机生成的8维“指纹”，同对象的重复测量只加入很小扰动；对象类别则独立随机生成，与指纹之间没有总体规律。一个最近邻规则在随机拆行时，可以从训练中找到几乎相同的对象记录；按对象拆分时则不能。

实际五折运行中，按行划分每折准确率都是1.0；按对象划分的五折准确率约为0.611、0.667、0.533、0.444、0.522。前者并非发现了随机类别的物理规律，而是利用了同对象的近重复记录。后者围绕机会水平波动，符合本人工构造没有跨对象预测信息的设置。

\fig{v3_ch04_groups.pdf}{同一人工数据的两种验证问题。按行拆分评估对已见对象的重复测量，按对象拆分评估新对象。两条曲线不能被解释为同一任务下两种算法的优劣。}
这个反例的结论不是“所有随机划分都有错”，而是\textbf{评价中的新样本必须与实际应用中的新样本含义一致}。若使用目标确实是预测已知对象的新测量，重复信息可能是允许的；若声称预测从未出现的对象，就必须相应隔离对象身份。

\section{分类与回归的评价指标}
\subsection{先数清楚哪些记录被分错}
\textbf{评价指标（metric）}是按明确规则把一批预测结果汇总为数值的方式。分类最直观的指标是\textbf{准确率（accuracy）}，即判断正确的记录数除以总数。为了知道错在哪一类，可以制作\textbf{混淆矩阵（confusion matrix）}：行表示参考类别，列表示预测类别，每个格子统计对应记录数。

考虑人工构造的10条二分类记录，把关注的类别称为正类，其余称为负类。“正”不是好、“负”不是坏，只是为了明确统计对象。表\ref{tab:confusion-hand}给出一组预测结果。
\begin{table}[htbp]\centering
\caption{用于手算的10条二分类记录汇总，不是实际观测结果。}\label{tab:confusion-hand}
\begin{tabular}{lrrr}\toprule
&预测为正&预测为负&合计\\\midrule
实际为正&3（TP）&1（FN）&4\\
实际为负&2（FP）&4（TN）&6\\\midrule
合计&5&5&10\\\bottomrule
\end{tabular}\end{table}

\textbf{真阳性TP}表示实际为正且预测为正；\textbf{假阳性FP}表示实际为负却报成正，也称误报；\textbf{假阴性FN}表示实际为正却判成负，也称漏检；\textbf{真阴性TN}表示实际为负且预测为负。先按实际与预测的两种状态定位格子，再记名称，就不容易混淆。

本例准确率为 $(3+4)/10=0.7$。如果关注预测为正的候选有多少可信，计算\textbf{精确率（precision）}：$3/(3+2)=0.6$。如果关注所有真正正类中找回多少，计算\textbf{召回率（recall）}：$3/(3+1)=0.75$。一般式为
\[
\mathrm{Precision}=\frac{\mathrm{TP}}{\mathrm{TP}+\mathrm{FP}},\qquad
\mathrm{Recall}=\frac{\mathrm{TP}}{\mathrm{TP}+\mathrm{FN}}.
\]
\textbf{F1分数}是二者的调和平均，在分母非零时为
\[
F_1=\frac{2\,\mathrm{Precision}\,\mathrm{Recall}}{\mathrm{Precision}+\mathrm{Recall}}
=\frac{2\mathrm{TP}}{2\mathrm{TP}+\mathrm{FP}+\mathrm{FN}}.
\]
本例为 $6/9=2/3$。分母为零时，应声明采用的软件约定，不能把未定义的情况当作有科学意义的零或满分。

\subsection{多分类与类别不平衡}
鸢尾花或葡萄酒有三个类别，混淆矩阵就是三行三列。逐类别计算召回率时，把当前类别视为正，其余合为负。\textbf{宏平均（macro average）}是先分别计算各类指标，再对类别等权平均；\textbf{加权平均}按各类样本数加权。类别很少的对象在两种汇总中的权重不同，报告时要说明选择。

例如1000条记录只有10条关注对象，一个全部预测为背景的规则准确率为99\%，却没有找回任何关注对象。这样的\textbf{类别不平衡}说明仅看准确率不够。指标应回答实际问题：漏掉对象的代价大，还是后续检查大量误报的代价大？
\fig{ch04_confusion.pdf}{鸢尾花的三分类混淆矩阵实例。对角线是判断正确的数量，非对角线给出具体的误判方向。}

\subsection{回归误差需要带着单位解释}
对回归，用 $e_i=\hat y_i-y_i$ 表示第 $i$ 条预测残差。\textbf{平均绝对误差MAE}先取绝对值再平均；\textbf{均方根误差RMSE}先平方取平均，再开平方：
\[
\mathrm{MAE}=\frac1n\sum_i|e_i|,\qquad
\mathrm{RMSE}=\sqrt{\frac1n\sum_i e_i^2}.
\]
假设人工给出目标 $(2,4,6)$ 与预测 $(3,4,5)$，残差为 $(1,0,-1)$。于是MAE为 $2/3\approx0.667$，RMSE为 $\sqrt{2/3}\approx0.816$。二者都具有目标的单位；平方再开根使较大的残差在RMSE中受到更明显的影响。

\textbf{决定系数 $R^2$}比较平方误差与以评价样本目标均值为常数的参照：
\[
R^2=1-\frac{\sum_i(y_i-\hat y_i)^2}{\sum_i(y_i-\bar y)^2}.
\]
上例 $\bar y=4$，分母为8，分子为2，所以 $R^2=0.75$。这不是“75\%的样本预测正确”。它可以为负，表示比该常数参照更差；目标完全不变时分母为零，需要另行处理。注意，此处评价样本均值只用于定义指标，真正可部署的常数预测器必须从训练目标估计常数，不能提前获取新样本目标\textsuperscript{[32]}。

\subsection{从混淆矩阵推到误报率和样本比例}
精确率的分母是所有被报为正类的记录，召回率的分母是真正的正类；\textbf{假阳性率}则是$\mathrm{FP}/(\mathrm{FP}+\mathrm{TN})$，分母是真正的负类。三者不能替换。对于极少见的对象，负类数目很大，即使假阳性率不高，也可能产生大量误报。

\worked{稀有对象筛选中的基数效应}
假设10万条记录中只有100条目标，模型召回90条目标，即召回率0.9。若假阳性率为0.01，则99900条背景中约有999条被误报，候选总数约1089，精确率只有$90/1089\approx0.0826$。这组数字是任务设定的算例，不是实际天文调查结果。

若模型在一个正负各半的评估集上表现很好，也不能直接把该集的精确率照搬到目标比例千分之一的巡天。预测概率的校准、类别比例以及后续核查预算都要重新考虑。这个例子解释了为什么候选筛选报告必须同时给出样本规模、类别比例、召回率和误报信息。

\subsection{阈值变化形成一组评价点}
概率分类器通常给出一个连续分数，把它转成类别需要阈值。逐渐降低阈值，报告的候选集合通常扩大；TP和FP都不会因集合扩大而减少，但精确率不保证单调。第6章将逐条计算一个阈值例子。ROC曲线以假阳性率为横轴、召回率为纵轴；PR曲线以召回率和精确率描述不同阈值的取舍。曲线概括排序行为，却不能替应用者决定最终阈值。

面积指标也不是一个应用阈值。例如较好的整体排序仍可能在高纯度要求的区域表现不足。对于科学候选目录，可以事先规定“人工最多核查50个对象”或“允许的误报预算”，然后在开发数据上选择满足目标的操作点，而不是仅凭曲线面积宣布模型足够好。

\subsection{同一个回归误差，可以有不同科学后果}
RMSE对大偏差更敏感，MAE较直接地报告典型绝对偏差。若误差单位是十万美元，数值0.4代表4万美元；若目标先取了对数，则对数误差需要转换后才能解释为原单位的相对偏差。不同目标变换后的分数不能直接排在同一表格里比较大小。

对于不确定度不同的观测，可以另外检查标准化残差$e_i/\sigma_i$，其中$\sigma_i$必须是明确来源的测量不确定度或预测尺度。它不是把误差“消除”，而是在相应噪声假设下比较偏离程度。\textbf{不能为使标准化残差看起来正常，事后任意调大分母。}本课程的主线先报告原单位误差，第16章再区分误差来源。

\worked{RMSE为什么不小于MAE}
对残差绝对值$a_i=|e_i|$，有
\[
\frac1n\sum_i a_i^2-\left(\frac1n\sum_i a_i\right)^2
=\frac1n\sum_i(a_i-\bar a)^2\ge0.
\]
两边开平方得到$\mathrm{RMSE}\ge\mathrm{MAE}$。只有所有残差绝对值相同时等号成立。这个短推导帮助检查程序：若同一批、同一权重的记录上报告RMSE小于MAE，应先核对实现、单位或数据范围，而不是急于解释模型。

\section{基线模型与公平比较}
\textbf{基线（baseline）}是用于比较的简单或已有方案。它提供“比什么更好”的参照。例如分类时始终报训练集中最多的类别，回归时始终报训练目标均值。它们未必足够好，却能检查一个复杂模型是否真的利用了输入信息。

假设训练分类记录中A占70\%，B占30\%，多数类基线对所有新输入都报A。新记录如果恰有同样比例，该基线也能有70\%准确率；所以得到71\%需要结合波动判断，不能单凭模型名称就称为显著进步。对于回归，训练目标为2、4、6时，均值基线给所有新输入报4。

sklearn提供 \texttt{DummyClassifier} 和 \texttt{DummyRegressor} 实现这类参照。\texttt{Dummy}在这里表示不利用复杂输入关系的简单规则，不是虚构数据。
\begin{lstlisting}[caption={先建立能解释清楚的多数类基线}]
from sklearn.dummy import DummyClassifier
from sklearn.model_selection import StratifiedKFold, cross_val_score
cv = StratifiedKFold(n_splits=5, shuffle=True, random_state=7)
baseline = DummyClassifier(strategy="most_frequent")
scores = cross_val_score(baseline, X_dev, y_dev, cv=cv, scoring="accuracy")
print(scores, scores.mean())
\end{lstlisting}

\textbf{公平比较}要求所比较方法面对相同的任务、样本范围、可用特征和评价划分。不同方法可以采用适合自身的预处理，但变换必须在各自训练范围内学习。若新模型同时增加了原模型没有的高质量测量，结果改善不能全部归因于算法。

\subsection{基线也要在每一折中重新拟合}
回归常数基线的均值来自当前训练目标，分类多数类来自当前训练标签。交叉验证中每折训练部分不同，基线参数也要重新确定。若预先使用全部数据的目标均值作为所有折的预测常数，验证信息已经影响了基线。虽然这种影响可能不大，完整比较仍应统一规则。

公平比较还需要限制可获得的信息和调参投入。例如模型A仅使用两个测量，模型B使用额外十个高精度测量，分数差异同时包含信息差异和算法差异。可以分别设置“同特征比较”和“增加信息比较”，而不是把二者混为一种进步。同理，一个模型搜索1000组配置，另一个只用默认配置，其结果既反映模型也反映开发投入。

\subsection{配对比较比只看两个平均数更有信息}
在完全相同的五折上评价模型A和B，可以计算每折差值$d_k=s_{A,k}-s_{B,k}$。若两模型分别在不同随机划分上得到平均0.95和0.96，无法判断差异主要来自模型还是验证样本难度。固定折至少可以让两者面对相同的难题。

但折间训练数据重叠，所以五个差值也不是五次完全独立的重复实验。对于本课程，先报告各折值、均值、差值方向和实际误判数量；不把微小平均差异自动写成统计显著。若需要正式比较，应依据对象相关结构和研究设计选择额外检验或重采样方法。

随书葡萄酒五折实验同时记录多数类基线和标准化逻辑回归。后者的作用此时只是一个提供好的候选规则，第6章再解释内部原理。不要因为它的名称更复杂而预先判断一定优于基线；先看同一折的实际结果。

\section{欠拟合、过拟合与模型复杂度}
\textbf{欠拟合（underfitting）}指当前模型或训练过程未充分表达数据中可学习的关系，常见表现是训练和验证误差都较大。\textbf{过拟合（overfitting）}指模型过度适应训练样本中的偶然细节，使训练改善没有转化为新样本上的改善。两者描述的是相对于任务、样本和训练过程的行为，不是固定的算法标签。

例如真实关系明显弯曲，模型却只允许一条直线，那么反复优化直线的斜率可能仍留下系统残差；这是表达能力不足的例子。如果允许一条曲线穿过每一个含噪测量点，它可能在训练点间剧烈弯折，对新测量反而预测更差；这是过度适应训练细节的例子。

\textbf{模型复杂度}描述规则表达多样关系或适应细节的能力，不能在所有算法间只用同一个数字衡量。后面的树模型通过增加条件分支形成更细的规则，神经网络通过增加计算单元改变表达能力。第3章的正则化则通过惩罚过大的参数施加限制。

图中树深可以先理解为沿一条判断路径允许进行多少层条件判断，具体树结构在第8章展开。曲线横轴改变复杂度，纵轴分别检查训练和验证结果；选择应以验证任务为依据，而不是只追求训练曲线最高。
\fig{ch04_tree_complexity.pdf}{条件判断层数增加时，训练与交叉验证表现可能不同。图中的树模型将在第8章解释，本章关注如何读两条评价曲线。}

\subsection{优化不足与表达不足不是同一问题}
训练误差高可能因为模型形式太简单，也可能因为参数还没有充分优化，还可能因为输入本来不含足够信息。三种情形的处理不同：增加训练轮数可能解决优化不足，却不能让一条直线变成弯曲函数；增加模型复杂度可能改善表达，却不能恢复未被测量的关键变量。

可以先对一个小训练子集检查模型是否能拟合，再查看训练目标的下降曲线和停止信息。如果小子集也学不好，先检查代码、数据配对、尺度和优化；如果训练可以非常好而验证差，才更有理由讨论过拟合和分布差异。这个诊断顺序在第12章神经网络中尤其重要，但也适用于其他模型。

\subsection{学习曲线和复杂度曲线}
复杂度曲线在固定数据量下改变模型配置；\textbf{学习曲线}在固定配置下改变训练样本量。两者横轴不同，不能混称。训练量增加时，训练分数可能下降，因为模型不能再只适应极少数记录；验证分数可能提高，也可能在模型形式受限时很快达到平台。

如果训练和验证误差都高且接近，继续收集相似样本未必解决表达不足；如果训练好、验证差且差距随样本增多而减小，增加数据可能有帮助。这些是诊断线索，不是仅凭曲线形状就能下的唯一结论。不同样本质量、划分方式和目标噪声也会改变曲线。

\subsection{选读：偏差与方差的一个精确分解}
假设$Y=f_*(x)+\varepsilon$，固定输入$x$，噪声满足均值零、方差$\sigma^2$，且与训练集独立。用随机训练集$D$得到预测$\hat f_D(x)$，记平均预测为$\bar f(x)=\mathbb E_D\hat f_D(x)$。将误差加减$\bar f$并展开，交叉项期望为零，得到
\[
\mathbb E_{D,\varepsilon}(Y-\hat f_D(x))^2
=\bigl(f_*(x)-\bar f(x)\bigr)^2
+\mathbb E_D\bigl(\hat f_D(x)-\bar f(x)\bigr)^2+\sigma^2.
\]
第一项是平均预测的偏离，第二项是训练数据变化导致的预测波动，第三项是这个输入信息下的随机噪声。这个等式需要上述假设，也只针对平方误差。它解释随机森林为什么关注降低预测波动，但不意味着在每个真实任务中可以从一套数据精确估计三项。

\section{交叉验证及结果波动}
\textbf{$K$折交叉验证（$K$-fold cross-validation）}将开发数据分为 $K$ 份，每次留一份验证，用其余份重新训练，轮流完成 $K$ 次。每一份称为一折。这里的 $K$ 是验证分组数，与第7章近邻算法的邻居数无关。

用12条人工编号记录说明三折验证：第一轮用5--12号训练、1--4号验证；第二轮用1--4和9--12号训练、5--8号验证；第三轮用1--8号训练、9--12号验证。每轮都重新建立模型与处理规则，不是让一个模型依次见到全部记录后再声称轮流验证。
\fig{v2_cross_validation.pdf}{12条人工编号记录的三折验证。每一行是一轮独立拟合，标出的验证部分在该轮不参与参数学习。}

若三轮准确率为0.75、1.00、0.75，平均值为 $2.5/3\approx0.833$。折间差异说明本次结果随划分变化，但这些训练部分相互重叠，折间标准差不是自动成立的置信区间。正式数据量与类别比例可能不同，也不应把三折手算数字写成葡萄酒实验结果。

\texttt{StratifiedKFold}尽量保持每折类别比例；\texttt{GroupKFold}按事先定义的对象组划分。一个材料有多次重复测量时，目标若是预测新材料，就不能把同一材料的重复测量同时放进训练与验证。只固定随机种子不能解决这种相关性问题。

\subsection{不相等的折大小与两种平均}
若一折验证20条、错2条，另一折验证30条、错6条，则两折准确率分别为0.9和0.8。简单折平均为0.85，把两折全部预测汇总后准确率为$(18+24)/50=0.84$。二者都可以计算，但含义不同。样本数几乎相等的普通$K$折中差别较小，按对象或时间分块时则可能明显。

对F1等非线性指标，先汇总混淆矩阵再计算，与先算各折指标再平均，也不一般相等。报告要写清采用哪种汇总。随书算法比较主要使用各折分数的算术平均，同时保存每折值；项目报告可以额外报告合并的折外预测指标，但应注明。

\subsection{嵌套验证：评价一个选择过程}
假设希望比较多个算法和参数，内层交叉验证可以负责选择；外层把一部分开发数据完全留出，用于检查内层选择流程的表现。每个外层训练部分都重新完成填补、缩放、调参和模型选择，再预测相应外层验证部分。这叫\textbf{嵌套交叉验证}。

它评价的是“按这套方法挑模型”这一过程，而不是某个事后挑出的固定参数。代价是拟合次数增加。小型课程练习可以采用固定测试集加内层验证，不强制每章嵌套；当报告涉及较多选择时，应理解为什么直接汇报最高内层分数容易乐观。

分组或时间约束要同时落实到内外两层。只在外层隔离天体，内层却按重复观测随机拆行，仍可能选出依赖对象重复信息的方案。验证设计是整个流程的属性，不是只改一个最后调用就能保证。

\section{超参数选择与模型选择}
第3章区分了参数与超参数。参数在训练中由数据学习，超参数是指定算法如何学习的配置，例如允许的树深、正则化强度或邻居数量。\textbf{模型选择}是根据验证结果在算法、特征处理和配置之间作出选择。

\textbf{网格搜索（grid search）}把少量候选值列成组合，逐个验证。例如配置A在三折得到 $(0.75,0.75,1.00)$，配置B得到 $(0.75,1.00,1.00)$，可以按事先规定的平均准确率偏向B；但这只是开发阶段选择，并不证明B在所有未来记录上都更好。最终保留数据用于提供额外检查。

后面的逻辑回归用 \texttt{C} 控制正则化相关配置，第6章解释它的作用。这里把它当作三个候选规则的配置标记即可，不需要提前掌握模型细节。\texttt{GridSearchCV}负责在开发数据上重复拟合并记录验证结果。
\begin{lstlisting}[caption={把候选配置限制在开发数据中比较}]
from sklearn.linear_model import LogisticRegression
from sklearn.preprocessing import StandardScaler
from sklearn.pipeline import Pipeline
from sklearn.model_selection import GridSearchCV
pipe = Pipeline([
    ("scale", StandardScaler()),
    ("clf", LogisticRegression(max_iter=2000))
])
search = GridSearchCV(pipe, {"clf__C": [0.1, 1.0, 10.0]},
                      cv=cv, scoring="accuracy")
search.fit(X_dev, y_dev)
print(search.best_params_, search.best_score_)
\end{lstlisting}

变量名 \texttt{clf\_\_C} 用双下划线连接步骤名和参数名。最高验证分数已经参与了方案选择，不能再把它标成未参与选择的测试分数。选择越多、越频繁地依据同一验证数据修改方案，越需要警惕对验证结果的偶然适应。

\subsection{从有限搜索到可解释的选择}
入门实验应该让每个候选值有理由。例如近邻数从1、3、5到较大的值，分别考察局部敏感与更强平滑；树深从1、2到更深，考察条件规则增加的效果。初次学习不宜把所有参数笛卡尔积组合成数千次搜索，因为学生很难解释每个变化。

选择准则也应事先确定。可以优先平均召回率，但要求精确率达到下限；可以优先RMSE，在差异很小的方案中选较简单或较便宜者。训练速度和预测成本同样可以记录，不过不同算法的运行条件必须一致。不能先看所有指标，再挑一个恰好支持自己偏好的结论。

\worked{选择后的分数为什么可能偏高}
假设十种实际能力相近的方案，在同一小验证集上的分数都包含随机波动。选择其中最高者，往往同时选中了正向波动较大的一次。这不需要某个算法有错误，也不需要故意作弊；只要“取最大值”，就可能产生选择偏差。另一个没有参与选择的数据集能够检查这个乐观程度。

提高实验可以人工生成完全随机的目标，增加候选模型数，观察最高验证分数怎样变化。该实验的意义是理解选择过程，不是说明真正数据中的所有高分都是假的。只有把目标设为随机、保持构造方式透明，才知道基准情形是什么。

\section{Pipeline与可复现的评价流程}
\subsection{流水线怎样避免漏掉处理步骤}
第2章已经手动分开了标准化的学习与应用。\textbf{Pipeline}将这些步骤与预测模型组成一个对象。在一次交叉验证中，程序先在当前训练折学习均值、尺度等规则，再训练模型；验证折只经过已经学习的变换和预测。下一折重新拟合整个流程，而不是沿用上一折的均值。

以上代码中的 \texttt{scale} 是标准化步骤的名称，\texttt{clf} 是分类模型步骤的名称。调用流水线 \texttt{fit} 时依次完成拟合与变换；调用 \texttt{predict} 时，对新数据应用同一预处理再预测。它不能自动阻止所有错误，例如一开始把目标列放进输入仍然会泄漏，但可以减少常见的预处理越界问题\textsuperscript{[3,4]}。

\subsection{什么叫可复现}
\textbf{可复现}在本书指另一位读者能依据记录的材料、软件和步骤，重新得到可核对的实验结果。至少保存数据来源、筛选条件、拆分索引、模型配置、软件版本、随机状态和评价代码。随机种子只固定部分随机过程，不保证跨版本完全相同，也不能代替稳定性检查。

完成一项预测实验后，应能回答：哪些记录用于学习参数，哪些用于选择方案，哪些直到最后才打开？某张图是训练结果还是验证结果？代码中的分数是单折、折均值还是最终保留结果？这些标注比一串小数位更重要。

\subsection{把实验记录写成可以复算的对象}
一个模型比较表至少包含：任务与目标单位、数据版本、开发与保留划分、特征列表、预处理、模型配置、指标及各折结果。不要只保存截图，因为截图通常不能恢复误判样本是谁。对应预测表应保留记录标识、参考目标、预测结果，以及必要的概率或残差。

流水线对象可以保存已经学习的处理规则，但不应替代清晰的代码和数据字典。模型文件由软件序列化产生，读取不可信来源的序列化对象可能执行代码，因此教学交付优先提供源码、数据和重训方式，不要求学生随意加载未知模型文件。

本书把“程序能够执行”“数值与手算一致”“验证设计回答正确问题”分开检查。第一项发现语法和依赖错误；第二项检查计算实现；第三项检查科学和统计设计。任何一项通过，都不能自动替代另外两项。这样的区分将在最后的项目评审中继续使用。

\chapterpractice{手算表\ref{tab:confusion-hand}的准确率、精确率、召回率和F1；计算三点回归例子的MAE、RMSE与$R^2$。用葡萄酒开发数据完成多数类基线与一个已提供候选模型的交叉验证，记录各折与均值，暂不反复查看最终测试。}
\chapterexercise
\begin{enumerate}
\item 把10条人工记录中一个假阳性改成真阴性，准确率、精确率和召回率分别怎样变化？
\item 训练、验证、测试三种用途能否仅凭文件名区分？举一个测试数据被用于选择方案的例子。
\item 对目标 $(2,4,6)$ 和预测 $(2,4,8)$，计算MAE与RMSE，并解释单位。
\item 写出12条记录三折验证的训练与验证编号。每轮为什么需要重新计算标准化参数？
\item 为什么在同一天体的重复观测之间随机拆行，不能直接代表对新天体的预测能力？
\end{enumerate}

\section*{深化练习与方法检验}
本组题强调解释、计算与验证；D编号在本章内独立计数。计算题应写出所用定义和中间步骤，实验题保留实际输出，包括不支持预期的结果。
\begin{enumerate}[label=D\arabic*.]
\item 两模型在同一100条验证记录上准确率分别为0.90和0.92。仅凭这两个数能否判断第二个模型的改善稳定？还应保存哪些逐记录信息，怎样设计进一步比较？
\item 真实正类占1\%，分类器召回率0.9、负类误报率0.05。以10000个对象计算预期真阳性、假阳性和精确率。解释为什么总体准确率不能代替候选清单的纯度。
\item 先在全部开发数据上选出与目标相关性最高的十个特征，再对这十列交叉验证，是否已经把特征选择纳入验证？写出正确的内外层用途。
\item 在本章重复对象模拟中，随机拆行的分数接近1，而按对象划分低得多。为什么不能简单说后者算法更差？怎样用部署问题决定应该报告哪种评价？
\end{enumerate}
